
\documentclass{article}
\usepackage{colortbl}
\usepackage{makecell}
\usepackage{multirow}
\usepackage{supertabular}

\begin{document}

\newcounter{utterance}

\centering \large Interaction Transcript for game `referencegame', experiment `line\_grids\_columns', episode 1 with reka{-}flash{-}3.1{-}t0.0.
\vspace{24pt}

{ \footnotesize  \setcounter{utterance}{1}
\setlength{\tabcolsep}{0pt}
\begin{supertabular}{c@{$\;$}|p{.15\linewidth}@{}p{.15\linewidth}p{.15\linewidth}p{.15\linewidth}p{.15\linewidth}p{.15\linewidth}}
    \# & $\;$A & \multicolumn{4}{c}{Game Master} & $\;\:$B\\
    \hline

    \theutterance \stepcounter{utterance}  
    & & \multicolumn{4}{p{0.6\linewidth}}{
        \cellcolor[rgb]{0.9,0.9,0.9}{
            \makecell[{{p{\linewidth}}}]{
                \texttt{\tiny{[P1$\langle$GM]}}
                \texttt{You are given three grids, where each of them is 5 by 5 in size.} \\
\texttt{Grids have empty cells marked with "▢" and filled cells marked with "X".} \\
\texttt{Your task is to generate a referring expression that best describes the target grid while distinguishing it from the two other distractor grids.} \\
\texttt{The first grid is the target grid, and the following two grids are the distractors.} \\
\\ 
\texttt{Target grid:} \\
\\ 
\texttt{□ □ X X X} \\
\texttt{□ □ X X X} \\
\texttt{□ □ X X X} \\
\texttt{□ □ X X X} \\
\texttt{□ □ X X X} \\
\\ 
\texttt{Distractor grid 1:} \\
\\ 
\texttt{X □ □ X X} \\
\texttt{X □ □ X X} \\
\texttt{X □ □ X X} \\
\texttt{X □ □ X X} \\
\texttt{X □ □ X X} \\
\\ 
\texttt{Distractor grid 2:} \\
\\ 
\texttt{□ X X X □} \\
\texttt{□ X X X □} \\
\texttt{□ X X X □} \\
\texttt{□ X X X □} \\
\texttt{□ X X X □} \\
\\ 
\texttt{Instruction: Describe the target grid.} \\
\texttt{Generate the referring expression starting with the tag "Expression: " for the given target grid. Omit any other text.} \\
            }
        }
    }
    & \\ \\

    \theutterance \stepcounter{utterance}  
    & \multicolumn{4}{p{0.6\linewidth}}{
        \cellcolor[rgb]{0.9,0.9,0.9}{
            \makecell[{{p{\linewidth}}}]{
                \texttt{\tiny{[P1$\rangle$GM]}}
                \texttt{<reasoning>} \\
\texttt{Let's tackle this problem. I need to generate a referring expression that describes the target grid and distinguishes it from the two distractor grids. I'll start by carefully looking at all three grids to understand their differences.} \\
\\ 
\texttt{First, the target grid. All rows are the same: five cells each, with the first two being empty (□) and the last three filled (X). So it's like two empty squares followed by three Xs in every row. That makes a solid block of Xs starting from the third column to the end, across all five rows. So the structure is consistent vertically and horizontally.} \\
\\ 
\texttt{Now, Distractor grid 1. Each row starts with an X, then three empty squares, and ends with an X. So the pattern here is X in the first and fifth columns, and empty in the middle three. Each row is like X □ □ □ X. So this forms a diagonal or maybe two Xs at the ends of each row. But since all rows have this, it's more like two Xs at the corners of each row, but not forming a diagonal across the grid. Wait, no, each row individually has X at first and last positions. So across the grid, it's like a border of Xs but only on the first and last columns. Wait, no, because each row has X in first and fifth, but the columns would have Xs in the first column (all rows) and fifth column (all rows). So the Xs form a vertical line on the left and right sides. So the grid has two vertical lines of Xs on the edges, and the middle three columns are empty. So the target grid is a horizontal block in the middle, while Distractor 1 is vertical lines on the sides.} \\
\\ 
\texttt{Distractor grid 2: Each row starts with an empty, then four Xs. Wait, looking at the description: "□ X X X □". Wait, no, the user wrote Distractor grid 2 as:} \\
\\ 
\texttt{□ X X X □} \\
\texttt{□ X X X □} \\
\texttt{□ X X X □} \\
\texttt{□ X X X □} \\
\texttt{□ X X X □} \\
\\ 
\texttt{Wait, each row has the first and last cells empty, and the middle three filled. Wait, no: the first cell is □, then X, X, X, then □} \\
            }
        }
    }
    & & \\ \\

    \theutterance \stepcounter{utterance}  
    & & & \multicolumn{2}{p{0.3\linewidth}}{
        \cellcolor[rgb]{0.9,0.9,0.9}{
            \makecell[{{p{\linewidth}}}]{
                \texttt{\tiny{[GM$|$GM]}}
                \texttt{Invalid generated expression} \\
            }
        }
    }
    & & \\ \\

\end{supertabular}
}

\end{document}
